\section{Primitive rings}
\label{Primitive rings}

We will consider (possibly non-unitary) rings. Thus  a \emph{ring} is an
abelian group $R$ with an associative multiplication $(x,y)\mapsto xy$ such
that $(x+y)z=xz+yz$ and $x(y+z)=xy+xz$ for all $x,y,z\in R$. If there is an
element $1\in R$ such that $x1=1x=x$ for all $x\in R$, we say that $R$ is a
\emph{unitary ring}.  A \emph{subring} $S$ of $R$ is an additive subgroup of
$R$ closed under multiplication. 

\begin{example}
    $\Z$ is a (unitary) ring and 
	$2\Z=\{2m:m\in\Z\}$ is a (non-unitary) ring.  
\end{example}

A \emph{left ideal} (resp. \emph{right ideal}) is a subring $I$ of $R$ such that 
$rI\subseteq I$ (resp. $Ir\subseteq I$) for all $r\in R$. An \emph{ideal}
(also two-sided ideal) of $R$ is a subring $I$ of $R$ that is both a left and a right ideal of $R$.

\begin{example}
	If $I$ and $J$ are both ideals of a ring $R$, then the sum 
 \[
 I+J=\{x+y:x\in I,y\in J\}
 \]
 and the intersection $I\cap J$ are both ideals of $R$. The product $IJ$,
 defined as the additive subgroup of $R$ generated by $\{xy:x\in I,y\in J\}$,
 is also an ideal of $R$. 
\end{example}

In the nonunital setting, $R^2=RR$ normally means the additive subgroup
generated by all products $xy$ for $x,y\in R$. 

\begin{example}
	If $R$ is a ring, the set $Ra =\{xa: x\in R\}$ is a left ideal
	of $R$. Similarly, the set $aR =\{ax: x\in R\}$ is a right ideal of $R$. The set $RaR$, which is
	defined as the additive subgroup of $R$ generated by $\{xay: x, y\in R\}$, is an 
	ideal of $R$.
\end{example}

\begin{example}
	If $R$ is a unitary ring, then $Ra$ is the left ideal generated by $a$, $aR$ is
	the right ideal generated by $a$ and $RaR$ is the ideal generated by $a$. 
	If $R$ is not unitary, the left ideal generated by $a$ is $Ra+\Z a$,
	the right ideal generated by $a$ is $aR+\Z a$ and the ideal generated by 
	$a$ is $RaR+Ra+aR+\Z a$.
\end{example}

The following exercise ask you to prove the \emph{Chinese Remainder Theorem}  
for arbitrary rings.

\begin{bonus}
    \label{xca:chinese}
    \index{Chinese Remainder Theorem}
    Let $R$ be a ring and $I_1,\dots,I_n$ be ideals such that 
    $I_j+I_k=R$ whenever $j\ne k$ and $R=I_j+R^2$ for all $j$. Prove that 
    \[
    R/(I_1\cap\cdots\cap I_n)\simeq R/I_1\times\cdots\times R/I_n.
    \]
\end{bonus}

In the previous exercise, the condition $R=I_j+R^2$ trivially holds in the case
of rings with identity. 

\begin{definition}
\index{Ring!simple}
A ring $R$ is said to be \emph{simple} if $R^2\ne\{0\}$ and the only ideals of 
$R$ are $\{0\}$ and~$R$.  
\end{definition}

The condition $R^2\ne\{0\}$ is trivially satisfied in the case of non-zero
rings with identity, as 
\[
1\in R^2=\{r_1r_2:r_1,r_2\in R\}.
\]

\begin{example}
	Division rings are simple.
\end{example}

Let $S$ be a unitary ring. Recall that for $n\geq1$, $M_n(S)$ is the ring of $n\times n$
square matrices with entries in $S$.  If $A=(a_{ij})\in M_n(S)$ and $E_{ij}$ is
the matrix such that $(E_{ij})_{kl}=\delta_{ik}\delta_{jl}$, then
\begin{equation}
	\label{eq:trick}
E_{ij}AE_{kl}=a_{jk}E_{il}
\end{equation}
for all $i,j,k,l\in\{1,\dots,n\}$. 

\begin{example}
	If $D$ is a division ring, then $M_n(D)$ is simple. 
\end{example}

Let $R$ be a ring. A left $R$-module (or module, for short)  
is an abelian group $M$ together with a map $R\times M\to M$, $(r,m)\mapsto r\cdot m$, such that
\begin{align*}
	&(r+s)\cdot m=r\cdot m+s\cdot m, &&
	r\cdot (m+n)=r\cdot m+r\cdot n, && r\cdot (s\cdot m)=(rs)\cdot m    
\end{align*}
for all $r,s\in R$, $m,n\in M$.  If $R$ has an identity 
$1$ and $1\cdot m=m$ holds for all $m\in M$, the module $M$ is said to be 
\emph{unitary}.  If $M$ is a unitary module, then $M=R\cdot M$. %\ne\{0\}$.

\begin{exercise}
\label{xca:center_simple}
Let $R$ be a simple unitary ring. 
\begin{enumerate}
    \item Prove that the center $Z(R)$ of $R$ is a field.
    \item Prove that $R$ is an algebra over $Z(R)$. 
\end{enumerate}
\end{exercise}

% Let $0\ne x\in Z(R)$. Then $Rx$ is a non-zero ideal of $R$.
% Since $R$ is simple, $Rx=R$. Thus $rx=1$ for some $r\in R$. 
% It follows that $x\in\mathcal{U}(R)$.

\begin{definition}
\index{Module!simple}
    A module $M$ is said to be 
    \emph{simple} if $R\cdot M\ne\{0\}$ and 
    the only submodules of $M$ are $\{0\}$ and $M$.
    If $M$ is a simple module, then $M\ne\{0\}$.
\end{definition}

If $R$ is a unitary ring and $M$ is a simple 
module, then $M$ is unitary. 

%\begin{remark}
%	Si $R$ es unitario y $M$ es un módulo simple, entonces $M$ es unitario.
%\end{remark}

\begin{lemma}
	\label{lemma:simple}
	Let $M$ be a non-zero module. Then $M$ is simple if and only if $M=R\cdot m$
	for all $0\ne m\in M$.
\end{lemma}

\begin{proof}
	Assume that $M$ is simple.  Let $m\ne 0$. Since $R\cdot m$ is a submodule of the simple 
	module $M$, either $R\cdot m=\{0\}$ or $R\cdot m=M$.  Let $N=\{n\in M:R\cdot n=\{0\}\}$. Since $N$ is a 
	submodule of $M$ and $R\cdot M\ne\{0\}$, $N=\{0\}$. Therefore $R\cdot m=M$, as $m\ne0$.
	Now assume that $M=R\cdot m$ for all $m\ne0$. Let $L$ be a non-zero submodule of 
	$M$ and let $0\ne x\in L$. Then $M=L$, as $M=R\cdot x\subseteq L$. 
\end{proof} 

\begin{example}
	Let $D$ be a division ring and let $V$ be a non-zero right vector space over $D$. If 
	$R=\End_D(V)$, then $V$ is a simple $R$-module with $fv=f(v)$ for
  $f\in R$ and 
	$v\in V$. 
% 	Para ver que $V$ es simple como $R$-módulo basta ver que $Rv=V$ para todo
% 	$v\ne0$. Sean $v,w\in V$, $v\ne0$.  Al completar $v\ne0$ a una base de $V$,
% 	vemos que existe $f\in R$ tal que $f(v)=w$. Luego $V$ es simple.
\end{example}

\begin{example}
	\label{exa:I_k}
	Let $n\geq2$.  If $D$ is a division ring and $R=M_n(D)$, then each 
	\[
	I_k=\{ (a_{ij})\in R:a_{ij}=0\text{ for $j\ne k$}\}
	\]
	is an $R$-module isomorphic to $D^n$. 
	Thus $M_{n}(D)$ is a simple ring that is not a simple $M_n(D)$-module.
\end{example}

\begin{definition}
\index{Minimal left ideal}
\label{def:minimal_left_ideal}
A left ideal $L$ of a ring $R$ is said to be \emph{minimal} if $L\ne\{0\}$ and 
$L$ does not properly contain non-zero left ideals of $R$. 
\end{definition}

Similarly one defines
minimal right ideals and minimal ideals. 

\begin{example}
	Let $n\geq1$, and let $D$ be a division ring and $R=M_n(D)$. Then $L=RE_{11}$ 
	is a minimal left ideal.
\end{example}

\begin{example}
	Let $L$ be a non-zero left ideal. If $RL\ne\{0\}$, then
	$L$ is minimal if and only if $L$ is a simple $R$-module. 
\end{example}

\begin{definition}
\index{Ideal!regular}
\index{Left ideal!regular}
\label{def:regular}
A left (resp. right) ideal $L$ of $R$ is said to be \emph{regular} if
there exists $e\in R$ such that $r-re\in L$ (resp.  $r-er\in L$) for all $r\in R$.
\end{definition}

If $R$ is a ring with identity, every left (or right) ideal is regular. 

\begin{definition}
\index{Ideal!maximal}
\index{Left ideal!maximal}
A left (resp. right) ideal $I$ of $R$ is said to be \emph{maximal} if $I\ne R$ and $I$ is not properly contained 
in a proper left (resp. right) ideal of $R$. 
\end{definition}

Similarly, one defines maximal ideals. 

A standard application of Zorn's lemma proves that every non-zero unitary ring
contains a maximal left ideal and a maximal right ideal. 

\begin{proposition}
	\label{proposition:R/I}
	Let $R$ be a ring and $M$ be a module. Then $M$ is simple if and only if
	$M\simeq R/I$ for some maximal regular left ideal $I$. 	
\end{proposition}

\begin{proof}
	Assume that $M$ is simple. Then $M=R\cdot m$ for some $m\ne0$ by 
	Lemma~\ref{lemma:simple}. The map $\phi\colon R\to M$, $r\mapsto r\cdot m$, 
	is a surjective homomorphism of $R$-modules, 
	so the first isomorphism theorem implies that 
	$M\simeq R/\ker\phi$. 
	
	We claim that $I=\ker\phi$ is a maximal left ideal. 
	The correspondence theorem 
	and the simplicity of $M$ imply that $I$ is a 
	maximal left ideal (because each left ideal $J$ such that 
	$I\subseteq J$ yields a submodule of $R/I$).

  We claim that $I$ is regular. Since $M=R\cdot m$, there exists $e\in R$ such
  that $m=e\cdot m$. If $r\in R$, then $r-re\in I$ since
  $\phi(r-re)=\phi(r)-\phi(re)=r\cdot m-r\cdot (e\cdot m)=0$.

    Now assume that $I$ is a maximal left ideal that is regular. 
    The correspondence theorem implies that 
    $R/I$ has no non-zero proper submodules. 
    
    We claim that 
    $R\cdot (R/I)\ne\{0_{R/I}\}$. Assume that $R\cdot (R/I)=\{0_{R/I}\}$ and let $r\in R$. 
    The regularity of $I$ implies that 
    there exists $e\in R$ such that $r-re\in I$. Hence $r\in I$, as  
	\[
    0_{R/I}=r\cdot (e+I)=re+I=r+I,
	\]
	a contradiction to the maximality of $I$. 
\end{proof}

\index{Annihilator!of a module}
\index{Annihilator!of a subset}
Let $R$ be a ring and $M$ be a left $R$-module. For a 
subset $N\subseteq M$
we define the \emph{annihilator} of $N$ as the subset 
\[
\Ann_R(N)=\{r\in R:r\cdot n=0\text{ for all }n\in N\}.
\]

\begin{example}
	$\Ann_{\Z}(\Z/n)=n\Z$.
\end{example}

\begin{exercise}
    Let $R$ be a ring and $M$ be a module. If $N\subseteq M$ is a subset, then
    $\Ann_R(N)$ is a left ideal of $R$. If $N\subseteq M$ is a submodule of
    $M$, then $\Ann_R(N)$ is an ideal of $R$. 
\end{exercise}

\begin{definition}
\index{Module!faithful}
A module $M$ is said to be \emph{faithful} if $\Ann_R(M)=\{0\}$. 
\end{definition}

\begin{example}
	If $K$ is a field, then $K^n$ is a faithful unitary $M_n(K)$-module.
\end{example}

\begin{example}
  If $V$ is a vector space over a field $K$, then $V$ is a faithful unitary
  $\End_K(V)$-module.
\end{example}

\begin{definition}
\index{Ring!primitive}
A ring $R$ is said to be \emph{primitive} if there exists a faithful simple $R$-module.  
\end{definition}

Since we work with left modules, ``primitive ideal'' means ``left primitive
ideal''. Right primitive ideals are defined similarly using simple right
modules.

\begin{exercise}
	\label{xca:simple=>prim}
	If $R$ is a simple unitary ring, then $R$ is primitive. 
\end{exercise}

\begin{exercise}
	\label{xca:prim+conm=cuerpo}
	If $R$ is a commutative ring (maybe without identity), then $R$ is primitive if and only if $R$ is a field. 
\end{exercise}

\begin{example}
	The ring $\Z$ is not primitive. 
\end{example}

\begin{definition}
\index{Ideal!primitive}
An ideal $P$ of a ring $R$ is said to be \emph{primitive} if $P=\Ann_R(M)$
for some simple $R$-module $M$. 
\end{definition}

As we work with left modules, our definition of primitive rings refers to
\emph{left primitive rings}. Of course, one can also define right primitive
rings. As Bergman showed, there exist rings that are right primitive but not
left primitive; see \cite{MR175940,MR167497}. 

\begin{exercise}
	\label{xca:primitivo}
	Let $R$ be a ring and $P$ be an ideal of $R$. Then $P$ is primitive if and only if 
	$R/P$ is a primitive ring.
\end{exercise}

% \begin{proof}
% 	Assume that $P=\Ann_R(M)$ for some $R$-module $M$. Then $M$ is a simple 
% 	$(R/P)$-module with 
% 	\[
% 	(r+P)\cdot m=r\cdot m,\quad r\in R,\;m\in M. 
% 	\]
% 	This operation 
% 	is well-defined, as 
% 	$P=\Ann_R(M)$. Since $M$ is a simple $R$-module, it follows that $M$ is 
% 	a simple $(R/P)$-module. Moreover, $\Ann_{R/P}M=\{0\}$. Indeed, if 
% 	$(r+P)\cdot M=\{0\}$, then $r\in\Ann_RM=P$ and hence $r+P=P$.

% 	Assume now that $R/P$ is primitive. Let $M$ be a faithful simple $(R/P)$-module. 
% 	Then 
%     \[
%     r\cdot m=(r+P)\cdot m, \quad r\in R,\; m\in M,
%     \]
%     turns $M$ into an $R$-module. It follows that $M$ is simple and that $P=\Ann_R(M)$. 
% \end{proof}

%\begin{example}
%	Si $I$ es un ideal maximal de un anillo unitario $R$, entonces $I$ es
%	primitivo. Como $I$ es ideal maximal y regular (pues $1\in R$), el cociente
%	$R/I$ es un anillo unitario simple y luego $R/I$ es primitivo por la
%	proposición~\ref{proposition:simple=>prim}. 
%\end{example}

%\begin{example}
%	Si $I$ es un ideal primitivo de un anillo conmutativo $R$, entonces $I$ es
%	maximal pues $R/I$ es un cuerpo (por ser primitivo y conmutativo), ver
%	proposición~\ref{proposition:prim+conm=cuerpo}.
%\end{example}

\begin{example}
	Let $n\geq1$ and $R_1,\dots,R_n$ be primitive rings and $R=R_1\times\cdots\times
	R_n$. Then each 
    \[
    P_i=R_1\times\cdots\times R_{i-1}\times\{0\}\times
	R_{i+1}\times\cdots\times R_n
    \]
    is a primitive ideal of $R$ since 
	$R/P_i\simeq R_i$.
\end{example}

%Recordemos que un ideal a izquierda $L$ de $R$ se dice \emph{minimal} si
%$L\ne0$ y $L$ no contiene propiamenete a otros ideales a izquierda no nulos de
%$R$.
%
%\begin{example}
%	Sea $L$ un ideal a izquierda de $R$ tal que $RL\ne0$. Entonoces $L$ es
%	simple si y sólo si $L$ es minimal.
%\end{example}

\begin{lemma}
	\label{lemma:maxprim}
	Let $R$ be a ring. If $P$ is a primitive ideal, there exists a regular 
        maximal left ideal $I$ such that $P=\{x\in R:xR\subseteq I\}$.
	Conversely, if $I$ is a regular maximal left ideal, then 
	$\{x\in R:xR\subseteq I\}$ is a primitive ideal. 
\end{lemma}

\begin{proof}
	Assume that $P=\Ann_R(M)$ for some simple $R$-module $M$. By
	Proposition~\ref{proposition:R/I}, there exists a regular maximal 
	left ideal 
	$I$ such that $M\simeq R/I$. Then 
    \[
    P=\Ann_R(R/I)=\{x\in
	R:xR\subseteq I\}.
    \]

	Conversely, let $I$ be a regular maximal left ideal. By
	Proposition~\ref{proposition:R/I}, $R/I$ is a simple $R$-module. Then
	\[
	\Ann_R(R/I)=\{x\in R:xR\subseteq I\}
	\]
	is a primitive ideal.
\end{proof}

%\begin{remark}
%	Una consecuencia trivial del lema~\ref{lemma:maxprim} es la siguiente: en
%	un anillo unitario, todo ideal a izquierda maximal contiene un ideal
%	primitivo.
%\end{remark}

\begin{exercise}
\label{xca:maximal=>primitive}
    Maximal ideals of unitary rings are primitive.  
\end{exercise}

\begin{exercise}
\label{xca:primitive=>maximal}
	Prove that every primitive ideal of a commutative ring is maximal.
\end{exercise}

\begin{bonus}
\label{xca:M_n(R)primitive}
  Let $n\geq1$. 
    Prove that $M_n(R)$ is primitive if and only if $R$ is primitive.
\end{bonus}

% Si $P$ es primitivo, entonces $R/P$ es un cuerpo(por ser primitivo y conmutativo) y luego $P$ es maximal
\section{Jacobson's radical}
\lecture\label{lec:Jacobson}

\begin{definition}
\index{Jacobson radical}
\label{def:J}
Let $R$ be a ring. The \emph{Jacobson radical} $J(R)$
is the intersection of all the annihilators of simple left $R$-modules. If $R$ does not
have simple left $R$-modules, then $J(R)=R$. 
\end{definition}

From the definition, it follows
that $J(R)$ is an ideal. Moreover, 
	\[
		J(R)=\bigcap\{P:\text{$P$ is a left primitive ideal of $R$}\}.
	\]

	If $I$ is an ideal of $R$ and $n\in\Z_{>0}$, $I^n$ is the additive subgroup of $R$ 
generated by the set $\{y_1\dots y_n:y_j\in I\}$. 

\begin{definition}
\index{Ideal!nilpotent}
\index{Left ideal!nilpotent}
An ideal $I$ of $R$ is \emph{nilpotent} 
if $I^n=\{0\}$ for some $n\in\Z_{>0}$.
\end{definition}

Similarly, one defines nilpotent left and right ideals. 
Note that an ideal $I$ is nilpotent if and only if there exists $n\in\Z_{>0}$ such that 
$x_1x_2\cdots x_n=0$ for all $x_1,\dots,x_n\in I$.  

\begin{definition}
    \index{Nil!element}
    \index{Nilpotent!element}
	An element $x$ of a ring is said to be \emph{nil} (or nilpotent) if $x^n=0$ for some $n\in\Z_{>0}$. 
\end{definition}

\begin{definition}
\index{Ideal!nil}
\index{Left ideal!nil}
An ideal $I$ of a ring is said to be \emph{nil} if every element of $I$ is nil. 
\end{definition}

Similarly, one defines nil left and right ideals. 
Note that every nilpotent ideal is nil, as $I^n=\{0\}$ implies $x^n=0$ for all 
$x\in I$.

\begin{example}
	Let \[
 R=\C[X_1,X_2,\dots]/(X_1,X_2^2,X_3^3,\dots), 
    \]
    and write $x_i$ for the image of $X_i$ in $R$. The ideal
    \[
    I=(x_1,x_2,x_3,\dots)
    \]
    is nil in $R$: every element of $I$ belongs to an ideal generated by
    finitely many commuting nilpotent elements, and is therefore nilpotent.
    However, $I$ is not nilpotent. Indeed, if $I$ is nilpotent, then 
    $I^k=\{0\}$ for some $k\geq1$. Thus $x_i^k=0$ for
    all $i$, a contradiction since $x_{k+1}^k\ne0$.
%    is nil in $R$, as it is generated by nilpotent elements. However, it is not nilpotent. Indeed, if $I$ is nilpotent, then there exists $k\in\Z_{>0}$ such that 
%	$I^k=\{0\}$ and hence $x_i^k=0$ for all $i$, a contradiction since 
%	$x_{k+1}^k\ne0$. 	
\end{example}

\begin{proposition}
	\label{pro:nilJ}
	Let $R$ be a ring. Then every nil left ideal is contained in $J(R)$.
\end{proposition}

\begin{proof}
	Assume that there is a nil left ideal $I$ such that 
	$I\not\subseteq J(R)$. There exists a simple $R$-module $M$ such that 
	$n=x\cdot m\ne 0$ for some $x\in I$ and $m\in M$. Since $M$ is simple,
	$R\cdot n=M$ and hence there exists $r\in R$ such that $r\cdot n=m$. Since
	\[
	(rx)\cdot m=r\cdot (x\cdot m)=r\cdot n=m, 
	\]
	it follows that $(rx)^k\cdot m=m$ 
    for all 
	$k\geq1$, a contradiction since $rx\in I$ is a nilpotent element. 
\end{proof}

Similarly, one proves that 
every nil right ideal is contained in the Jacobson radical. 


\begin{definition}
\index{Element!left quasi-regular}
\index{Element!right quasi-regular}
\index{Element!quasi-regular}
Let $R$ be a ring. An element $a\in R$ is said to be 
\emph{left quasi-regular} if there exists $r\in R$ such that $r+a+ra=0$. Similarly, 
$a$ is said to be \emph{right quasi-regular} if there exists $r\in R$ such that $a+r+ar=0$. An element is \emph{quasi-regular} if it is both left and right
quasi-regular.
\end{definition}

% \begin{exercise}
% 	\label{exercise:circ}
Let $R$ be a ring. A direct calculation shows that
the \emph{Jacobson circle operation}\index{Jacobson circle operation}
\[
R\times R\to R,
\quad
(r,s)\mapsto r\circ s=r+s+rs,
\]
is an associative operation with neutral element $0$.

\begin{example}
Let $R=\Z/3=\{0,1,2\}$ be the ring of integers modulo 3. 
The Jacobson circle 
operation of $R$ is shown in Table~\ref{tab:radical}.
\end{example}

	\begin{table}[bht]
  		\caption{The table of the Jacobson circle operation on $\Z/3$.}
		\centering
		\begin{tabular}{c|ccc}
			$\circ$ & 0 & 1 & 2\tabularnewline
			\hline
			0 & 0 & 1 & 2\tabularnewline
			1 & 1 & 0 & 2\tabularnewline
			2 & 2 & 2 & 2\tabularnewline
		\end{tabular}
    \label{tab:radical}
 	\end{table}

%\end{exercise}

% \begin{exercise}
% 	Let $R=\Z/3=\{0,1,2\}$. Compute the multiplication table with respect to the circle 
%  	operation given by the previous exercise.  
%  	%is then 
% % 	\begin{table}[ht]
% % 		\centering
% % 		\begin{tabular}{c|ccc}
% % 			$\circ$ & 0 & 1 & 2\tabularnewline
% % 			\hline
% % 			0 & 0 & 1 & 2\tabularnewline
% % 			1 & 1 & 0 & 2\tabularnewline
% % 			2 & 2 & 2 & 2\tabularnewline
% % 		\end{tabular}
% % 		\caption{The multiplication table of the radical ring $\Z/3$.}
% % 	\end{table}
% \end{exercise}

If $R$ is unitary, an element $x\in R$ is left quasi-regular (resp. right
quasi-regular) if and only if $1+x$ is left invertible (resp. right
invertible). In fact, if $r\in R$ is such that $r+x+rx=0$, then
$(1+r)(1+x)=1+r+x+rx=1$.  Conversely, if there exists $y\in R$ such that
$y(1+x)=1$, then  
\[
(y-1)\circ x=y-1+x+(y-1)x=0.
\]

\begin{example}
	If $x\in R$ is a nilpotent element,
    then $y=\sum_{n\geq1}x^n\in R$ is left quasi-regular.
	In fact, if there exists $N$ such that $x^N=0$,
    then the sum defining $y$ is finite and
    \[
    (-x)+y+(-x)y=0.
    \]
    Is $y$ also right quasi-regular?
%	If $x\in R$ is a nilpotent element, 
%    then $y=\sum_{n\geq1}x^n\in R$ is left quasi-regular. 
%	In fact, if there exists $N$ such that $x^N=0$, 
%    then the sum defining $y$ is finite 
%    and $y+(-x)+y(-x)=0$.  Is it also right quasi-regular?
%	En efecto, si existe $N$ tal que $x^N=0$, la suma que
%	define al elemento $y$ es finita y cumple que $(-x)+y+(-x)y=0$.  
\end{example}

\begin{definition}
    \index{Left ideal!quasi-regular}
A left ideal $I$ of $R$ is said to be 
\emph{left quasi-regular} (resp. right quasi-regular) if every element of $I$ is
left quasi-regular (resp. right quasi-regular). A left ideal 
is said to be \emph{quasi-regular} if it is left and right quasi-regular. 
\end{definition}

Similarly 
one defines right quasi-regular ideals and quasi-regular ideals. 

\begin{lemma}
	\label{lemma:casiregular}
	Let $I$ be a left ideal of $R$. If $I$ is left quasi-regular, then 
	$I$ is quasi-regular.
\end{lemma}

\begin{proof}
	Let $x\in I$. Let us prove that $x$ is right quasi-regular. Since $I$ is
	left quasi-regular, there exists $r\in R$ such that $r\circ x=r+x+rx=0$. Since 
	$r=-x-rx\in I$, there exists $s\in R$ such that $s\circ
	r=s+r+sr=0$. Then $s$ is right quasi-regular and  
	\[
	x=0\circ x=(s\circ r)\circ x=s\circ (r\circ x)=s\circ 0=s.\qedhere
	\]
\end{proof}

% \index{Lemma!Zorn}
% Let $(A,\leq)$ be a \emph{partially order set}, this means that $A$ is a set together with a 
% reflexive, transitive, and anti-symmetric binary relation
% $R$ en $A\times A$, where $a\leq b$ if and only if $(a,b)\in R$. 
% Recall that the relation is reflexive if $a\leq a$ for all $a\in A$, the relation is transitive if 
% $a\leq b$ and $b\leq c$ imply that 
% $a\leq c$ and the relation is anti-symmetric if $a\leq b$ and $b\leq a$ imply $a=b$.
% The elements $a,b\in A$ are said to be \emph{comparable} if $a\leq b$ or $b\leq
% a$. An element $a\in A$ is said to be \emph{maximal} if 
% $c\leq a$ 
% for all $c\in A$
% that is comparable with $a$. 
% An \emph{upper bound} for a non-empty subset $B\subseteq A$ is an element $d\in
% A$ such that $b\leq d$ for all $b\in B$. A \emph{chain} in $A$ is a subset 
% $B$ such that every pair of elements of $B$ are comparable. 
% \emph{Zorn's lemma} states the following property: 
% \begin{quote}
% If $A$ is a non-empty partially ordered set such that every chain in 
% $A$ contains an upper bound in $A$, then $A$ contains a maximal element. 
% \end{quote}

% Our application of Zorn's lemma:
For unitary rings, every proper left ideal is contained in a maximal left
ideal; this is a standard application of Zorn's lemma. A related result for
arbitrary rings is the following.

\begin{lemma}
	\label{lemma:maxreg}
	Let $R$ be a ring, and $x\in R$ be an element that is not left quasi-regular. Then there
	exists a maximal left ideal $M$ such that 
	$x\not\in M$. Moreover, $R/M$ is a simple $R$-module and 
	$x\not\in\Ann_R(R/M)$.
\end{lemma}

\begin{proof}
  Let $T=\{r+rx:r\in R\}$. A straightforward calculation shows that $T$ is a
  left ideal of $R$ such that $x\not\in T$. Indeed, if $x\in T$, then $-x\in
  T$, so $r+rx=-x$ for some $r\in R$, contradicting the assumption that $x$ is
  not left quasi-regular. 

  The only left ideal of $R$ containing $T\cup\{x\}$ is $R$. Indeed, if $U$ is
  a left ideal containing $T\cup\{x\}$, then every $r\in R$ can be written as
%	$R$ such that $x\not\in T$ (if $x\in T$, then $r+rx=-x$ for some 
%	$r\in R$, a contradiction since $x$ is not left quasi-regular). 
%	The only left ideal of $R$ containing 
%	$T\cup\{x\}$ is $R$. Indeed, if there exists a left ideal $U$ containing $T$, then 
%    $x\not\in U$, since otherwise every $r\in R$ could be written as 
	\[
    r=(r+rx)+r(-x)\in U,
    \]
  and hence $U=R$.

  Let $\mathcal{S}$ be the set of proper left ideals of $R$ containing $T$
  partially ordered by inclusion. Since $T\in\mathcal{S}$,
  $\mathcal{S}\ne\emptyset$.  If $\{K_i:i\in I\}$ is a chain in $\mathcal{S}$,
  then $K=\cup_{i\in I}K_i$ is an upper bound for the chain 
	($K$ is proper, as $x\not\in K$). Then Zorn's lemma implies that 
	$\mathcal{S}$ admits a maximal element $M$. Thus $M$
	is a maximal left ideal such that $x\not\in M$.
 
  Moreover, $M$ is regular
%	since $r-r(-x)\in T\subseteq M$ for all $r\in R$. Therefore $R/M$ is a simple 
%	$R$-module by Proposition~\ref{proposition:R/I}. Since $x\cdot (x+M)\ne
%	M$ (if $x^2\in M$, then  $x\in M$, as $x+x^2\in
	since $r-r(-x)\in T\subseteq M$ for all $r\in R$. Therefore $R/M$ is a simple
	$R$-module by Proposition~\ref{proposition:R/I}. Since
    $x\cdot(x+M)=x^2+M\ne M$ (if $x^2\in M$, then $x\in M$, as
    $x+x^2\in
	T\subseteq M$), it follows that $x\not\in\Ann_R(R/M)$.
\end{proof}

\begin{exercise}
\label{xca:not_in_J(R)}
Let $R$ be a ring and 
$x\in R$ be an element that is not left quasi-regular. Prove that $x\not\in J(R)$. 
%there exists 
%, the lemma implies that there exists 
%a simple $R$-module $M$ such $x\not\in\Ann_R(M)$. Thus 
%$x\not\in J(R)$.
\end{exercise}

\begin{theorem}
	\label{thm:casireg_eq}
	Let $R$ be a ring and $x\in R$. The following statements are equivalent: 
	\begin{enumerate}
		\item The left ideal generated by $x$ is quasi-regular.
		\item $Rx$ is quasi-regular.
		\item $x\in J(R)$.
	\end{enumerate}
\end{theorem}

\begin{proof}
	The implication $(1)\implies(2)$ is trivial, as $Rx$ is included in the left ideal 
	generated by $x$.  
	
	We now prove $(2)\implies(3)$. If
	$x\not\in J(R)$, by definition, there exists a simple 
	$R$-module $M$ such that $x\cdot m\ne 0$ for some $m\in M$. The simplicity of $M$ implies
	that $(Rx)\cdot m=M$. Thus there exists $r\in R$ such that $(rx)\cdot m=-m$. There is an element 
	$s\in R$ such that $s+rx+s(rx)=0$ and hence 
	\[
	-m=(rx)\cdot m=(-s-srx)\cdot m=-s\cdot m+s\cdot m=0,
	\]
	a contradiction. 
	
	Finally, we prove $(3)\implies(1)$. Let $x\in J(R)$ and $I$ be the left ideal of $R$ 
 generated by $x$. Since $I\subseteq J(R)$ and every element of $J(R)$ is left quasi-regular (see Exercise~\ref{xca:not_in_J(R)}), 
 $I$ is left quasi-regular. By Lemma~\ref{lemma:casiregular}, $I$ is quasi-regular. 
\end{proof}

The theorem implies the following corollary. 

\begin{corollary}
	If $R$ is a ring, then $J(R)$ is a quasi-regular ideal that contains every 
	quasi-regular left ideal. 
\end{corollary}

\begin{proof}
    By Exercise~\ref{xca:not_in_J(R)}, $J(R)$ is left quasi-regular. By Lemma~\ref{lemma:casiregular}, $J(R)$ is quasi-regular. 
    Let $I$ be a left ideal of $R$ that is quasi-regular. Let $x\in I$ and $J$ be the left
    ideal of $R$ generated by $x$. Since $J\subseteq I$, $J$ is quasi-regular. By Theorem~\ref{thm:casireg_eq}, $x\in J(R)$. 
\end{proof}

% The following result is somewhat what we all had in mind. We first need a lemma. 

% \begin{lemma}
%     \label{lemma:JsupsetCR}
%     Let $R$ be such that $J(R)\ne R$. If $I$ is a left quasi-regular left 
%     ideal of $R$, then $I\subseteq J(R)$.
% \end{lemma}

% \begin{proof}
%     Assume that $I\not\subseteq J(R)$. There exists a simple $R$-module 
%     $N$ such that $I\cdot N\ne \{0\}$. In particular, $I\cdot n\ne \{0\}$ for some $0\ne n\in N$. Since 
%     $I$ is a left ideal, $I\cdot n$ is a non-zero submodule of $N$. 
%     Then $I\cdot n=N$, as $N$ is simple. There exists $x\in I$ such that $x\cdot n=-n$. 
%     Since $I$ is left quasi-regular, there exists $r\in R$ such that $r+x+rx=0$.
%     Thus 
%     \[
% 	0=0\cdot n=(r+x+rx)\cdot n=r\cdot n+x\cdot n+(rx)\cdot n=r\cdot n-n-r\cdot n=-n,
%     \]
%     a contradiction. 
% \end{proof}


% \begin{proof}
% 	Es consecuencia inmediata del teorema~\ref{thm:casireg_eq}. 
%%	Como \[
%%	J(R)=\bigcap\{\Ann_R(R/I):I\text{ ideal a izquierda maximal y
%%	regular}\}
%%	\]
%%	por el teorema~\ref{thm:J(R)}, 
%%	todo $x\in J(R)$ es casi-regular gracias al
%%	lema~\ref{lemma:K}. Si $I$ es un ideal casi-regular a izquierda, entonces
%%	$I\subseteq J(R)$ por el lema~\ref{lemma:JsupsetCR}.
%\end{proof}

%\begin{lemma}
%%	\label{lemma:maxreg}
%	Sea $R$ un anillo y sea $I\ne R$ un ideal a izquierda regular. Entonces $I$
%	está contenido en algún ideal a izquierda maximal y regular.
%\end{lemma}
%
%\begin{proof}
%	
%\end{proof}

%\begin{lemma}
%	\label{lemma:K}
%	Sean $R$ un anillo y $K=\bigcap\{I:\text{$I$ ideal a izquierda maximal y
%	regular}\}$.  Entonces $K$ es un ideal a izquierda casi-regular.
%\end{lemma}
%
%\begin{proof}
%	Gracias al lema~\ref{lemma:casiregular}, basta ver que $K$ es casi-regular
%	a izquierda.  Sea $a\in K$ y sea $T=\{r+ra:r\in R\}$.  Es claro que $T$ es
%	un ideal a izquierda regular con $e=-a$. Como $a\in K$, existe $s\in R$ tal
%	que $s+a+sa=0$. Entonces, como $-a=s+sa$, 
%	\[
%	r+re=r+r(-a)=r+r(s+sa)=r+rs+rsa\in r+T
%	\]
%	para todo $r\in R$. 
%	Si $T\ne R$, por el lema~\ref{lemma:maxreg} existe un ideal a izquierda
%	$J$, maximal y regular.  Como $a\in K\subseteq J$, $J$ es ideal a izquierda
%	y $r+ra\in T\subseteq J$ para todo $r\in R$, se concluye que $R=J$, una
%	contradicción.  Luego $T=R$ y entonces existe $r\in R$ tal que $r+ra=-a$.
%	Esto implica que $a$ es casi-regular a izquierda. 
%\end{proof}

%\begin{lemma}
%	\label{lemma:JsupsetCR}
%	Sea $R$ un anillo que admite un $R$-módulo simple. Si $I$ es un ideal a
%	izquierda casi-regular a izquierda, $I\subseteq J(R)$.
%\end{lemma}
%
%\begin{proof}
%	Supongamos que $I\not\subseteq J(R)$. Existe entonces un $R$-módulo simple
%	$N$ tal que $IN\ne 0$. Entonces $In\ne 0$ para algún $0\ne n\in N$. Como
%	$I$ es un ideal a izquierda, $In\subseteq N$ es un submódulo no nulo del
%	simple $N$. Luego $In=N$. Existe entonces $x\in I$ tal que $xn=-n$. Como
%	$I$ es casi-regular a izquierda, existe $r\in R$ tal que $r+x+rx=0$.
%	Entonces
%	\[
%		0=0n=(r+x+rx)n=rn+xn+rxn=rn-n-rn=-n,
%	\]
%	una contradicción.
%\end{proof}

The following exercise uses Zorn's lemma 
and will be used in the proof of Theorem~\ref{thm:J(R)}. 

\begin{exercise}
\label{xca:regular}
    Let $R$ be a ring. 
    Prove that every proper left ideal of $R$ that 
    is regular is contained in a maximal left ideal that is regular. 
\end{exercise}

\begin{theorem}
	\label{thm:J(R)}
	Let $R$ be a ring such that $J(R)\ne R$. Then 
	\begin{align*}
		J(R)&=\bigcap\{I:\text{$I$ is a regular maximal left ideal of $R$}\}.
	\end{align*}
\end{theorem}

\begin{proof}
        Let 
	\[
	K=\bigcap\{I:\text{$I$ is a regular maximal left ideal of $R$}\}.
	\]
        
        Let us prove that $K\subseteq J(R)$. By the preceding corollary, it is
        enough to prove that $K$ is quasi-regular. By
        Lemma~\ref{lemma:casiregular}, it is enough to prove that $K$ is left
        quasi-regular. 

        Let $a\in K$ and $T=\{r+ra:r\in R\}$. 
        If $T=R$, then $-a=r+ra$ for some $r\in R$ and hence 
        $a$ is left quasi-regular. So we need to prove that $T=R$. Note that 
        $T$ is a regular left ideal with $e=-a$ (see
        Definition~\ref{def:regular}). If $T\ne R$, then, by the previous
        exercise, $T$ is contained in a regular maximal left ideal $M$. Then
        $a\in K\subseteq M$ and hence $ra\in M$ for all $r\in R$. Since
        $r+ra\in T\subseteq M$ for all $r\in R$, it follows that $M=R$, a
        contradiction. Therefore $T=R$. 
%        $T$ is a regular left ideal with $e=-a$ (see Definition~\ref{def:regular}). If $T\ne R$, then $T$ is contained in a maximal left ideal $J$  
%        by the previous exercise. Then 
%        $a\in K\subseteq J$ and hence $ra\in J$ for all $r\in R$. Since 
%        $r+ra\in T\subseteq J$ for all $r\in R$, it follows that $J=R$, a contradiction. Therefore $T=R$. 
        
        Now we prove that $J(R)\subseteq K$. 
		By
	Proposition~\ref{proposition:R/I}, 
	\[
		J(R)=\bigcap\{\Ann_R(R/I):I\text{ is a regular maximal left ideal of $R$}\}.
	\]
	Let $I$ be a regular maximal left ideal. If $r\in J(R)\subseteq
	\Ann_R(R/I)$, then, since $I$ is regular, there exists $e\in R$ such that
	$r-re\in I$. Since 
    $r\in\Ann_R(R/I)$, 
	\[
	%re+I=r(e+I)\subseteq I,
    r\cdot(e+I)=re+I=I.
	\]
	Thus $re\in I$ and hence $r\in I$. Therefore  $J(R)\subseteq K$. 
\end{proof}

\begin{example}
  Every maximal ideal of $\Z$ is of the form $p\Z=\{pm:m\in\Z\}$ for some
  prime number $p$.  Thus $J(\Z)=\bigcap_p p\Z=\{0\}$.
\end{example}

%\begin{example}
%	Sea $D$ un anillo de división y sea $R=D[x_1,\dots,x_n]$. Si $f\in J(R)$
%	entonces\dots Luego $J(R)=0$. 
%\end{example}

We now review some basic results useful to compute radicals. 

\begin{proposition}
	Let $\{R_i:i\in I\}$ be a family of rings. Then 
	\[
	J\left(\prod_{i\in I}R_i\right)=\prod_{i\in I}J(R_i).
	\]
\end{proposition}

\begin{proof}
  Let $R=\prod_{i\in I}R_i$ and $x=(x_i)_{i\in I}\in R$. Then
  \[
    Rx=\prod_{i\in I}R_ix_i.
  \]
  An element of $R$ is quasi-regular if and only if each of its coordinates is
  quasi-regular. Therefore $Rx$ is quasi-regular if and only if every $R_ix_i$
  is quasi-regular in $R_i$. By Theorem~\ref{thm:casireg_eq}, $x\in J(R)$ if
  and only if $x_i\in J(R_i)$ for all $i\in I$.
%	Let $R=\prod_{i\in I}R_i$ and $x=(x_i)_{i\in I}\in R$.  The left ideal 
%    $Rx$ is quasi-regular if and only if each left ideal $R_ix_i$
%	is quasi-regular in $R_i$, as $x$ is quasi-regular in $R$ if and only if each 
%	$x_i$ is quasi-regular in $R_i$. Thus $x\in J(R)$ if and only if $x_i\in
%	J(R_i)$ for all $i\in I$.
\end{proof}

For the next result, we shall need a lemma.

\begin{lemma}
	\label{lemma:trickJ1}
	Let $R$ be a ring and $x\in R$. 
	If $-x^2$ is a left quasi-regular element, then so is $x$. 
\end{lemma}

\begin{proof}
	Let $r\in R$ be such that $r+(-x^2)+r(-x^2)=0$ and $s=r-x-rx$. Then
    $x$ is left quasi-regular, as 
    \begin{align*}
		s+x+sx&=(r-x-rx)+x+(r-x-rx)x\\
		&=r-x-rx+x+rx-x^2-rx^2=r-x^2-rx^2=0.\qedhere 
\end{align*}
\end{proof}

%\begin{lemma}
%	\label{lemma:trickJ2}
%	Sea $R$ un anillo. Entonces $x\in J(R)$ si y sólo si $Rx$ es un ideal a
%	izquierda casi-regular a izquierda.
%\end{lemma}
%
%\begin{proof}
%	Si $x\in J(R)$ entonces $Rx\subseteq J(R)$ y luego todo elemento de $Rx$ es
%	casi-regular a izquierda. Recíprocamente, si $Rx$ es casi-regular a
%	izquierda, $Rx+\Z x$ es un ideal a izquierda de $R$. Si $s=rx+nx\in Rx+\Z
%	x$, entonces $-s^2$ es casi-regular a izquierda (pues $-s^2\in Rx$). Por el
%	lema~\ref{lemma:trickJ1}, $s$ es casi-regular a izquierda; en particular,
%	$x$ es casi-regular a izquierda y luego $x\in J(R)$. 
%\end{proof}

\begin{proposition}
	\label{proposition:J(I)}
	If $I$ is an ideal of $R$, then $J(I)=I\cap J(R)$. 
\end{proposition}

\begin{proof}
  Note that $I\cap J(R)$ is an ideal of $I$. Let $x\in I\cap J(R)$ and $a\in
  I$. Since $ax\in J(R)$, the element $ax$ is left quasi-regular in $R$. Thus
  there exists $s\in R$ such that $s+ax+sax=0$. Since $s=-ax-sax\in I$, the
  element $ax$ is left quasi-regular in $I$.  Therefore $Ix$ is a left
  quasi-regular left ideal of $I$; by Lemma~\ref{lemma:casiregular} it is
  quasi-regular. Applying Theorem~\ref{thm:casireg_eq} to the ring $I$ gives
  $x\in J(I)$. Hence $I\cap J(R)\subseteq J(I)$.

  Conversely, let $x\in J(I)\subseteq I$ and $r\in R$. Since
  \[
    -(rx)^2=(-rxr)x\in I\,J(I)\subseteq J(I),
  \]
  the element $-(rx)^2$ is left quasi-regular in $I$, and hence in $R$.
  Lemma~\ref{lemma:trickJ1}, applied in $R$, shows that $rx$ is left
  quasi-regular. Thus $Rx$ is a left quasi-regular left ideal of $R$, so it is
  quasi-regular by Lemma~\ref{lemma:casiregular}. By
  Theorem~\ref{thm:casireg_eq}, $x\in J(R)$. Therefore $J(I)\subseteq I\cap
  J(R)$.
\end{proof}

\begin{optional}
\section{Radical rings} 

\begin{definition}
\index{Ring!radical}
A ring $R$ is said to be \emph{radical} if $J(R)=R$. 
\end{definition}

\begin{example}
	If $R$ is a ring, then $J(R)$ is a radical ring, by Proposition~\ref{proposition:J(I)}.
\end{example}

\begin{example}
	The Jacobson radical of $\Z/8$ is $\{0,2,4,6\}$. 
\end{example}

There are several characterizations of radical rings. 

\begin{theorem}
	\label{theorem:anillo_radical}
	Let $R$ be a ring. The following statements are equivalent: 
	\begin{enumerate}
		\item $R$ is radical.
		\item $R$ admits no simple $R$-modules. 
		\item $R$ does not have regular maximal left ideals.
		\item $R$ does not have left primitive ideals.
		\item Every element of $R$ is quasi-regular. 
		\item $(R,\circ)$ is a group. 
	\end{enumerate}
\end{theorem}

\begin{xca}
    Prove Theorem~\ref{theorem:anillo_radical}. 
\end{xca}
%\begin{proof}
%	The equivalence $(1)\Longleftrightarrow(5)$ follows from 
%	Theorem~\ref{thm:casireg_eq}. 
%    
%    The equivalence $(5)\Longleftrightarrow(6)$ is left as an exercise. 
%
%	Let us prove that $(1)\implies(2)$. Assume that there exists a simple $R$-module $N$. Since 
%	$R=J(R)\subseteq\Ann_R(N)$, $R=\Ann_R(N)$. 
%	Hence $R\cdot N=\{0\}$, a contradiction to the simplicity of $N$.
%	
%	To prove $(2)\implies(3)$ we note that for each regular and maximal left ideal 
%	$I$, the quotient $R/I$ is a simple $R$-module by
%	Proposition~\ref{proposition:R/I}. 
%	
%	To prove $(1)\implies(4)$ assume that there is a primitive left ideal 
%	$I=\Ann_R(M)$, where $M$ is some simple $R$-module. Since $R=J(R)\subseteq I$, it follows that  
%    $I=R$, a contradiction to the simplicity of $M$.
%
%	Finally we prove $(4)\implies(2)$. If $M$ is a simple $R$-module, then 
%	$\Ann_R(M)$ is a primitive left ideal.
%\end{proof}

\begin{example}
	Let 
	\[
	A=\left\{\frac{2x}{2y+1}:x,y\in\Z\right\}.
	\]
	Then $A$ is a radical ring, as the inverse of an arbitrary element of the form $\frac{2x}{2y+1}$
	with respect to the Jacobson circle operation 
	is 
    \[
	\left(\frac{2x}{2y+1}\right)'=\frac{-2x}{2(x+y)+1}.
	\]
\end{example}

\begin{definition}
A pair $(X,r)$ is a \emph{solution} to the Yang--Baxter equation (YBE) if $X$
is a non-empty set and $r\colon X\times X\to X\times X$ is a bijective map such
that  
\[
	(r\times\id)\circ (\id\times r)\circ (r\times\id)
	=(\id\times r)\circ (r\times\id)\circ (\id\times r).
\]
\end{definition}

If $(X,r)$ is a solution (to the YBE)\index{Yang--Baxter equation}, 
by convention we write 
\[
	r(x,y)=(\sigma_x(y),\tau_y(x)).
\]

We say that a solution $(X,r)$ is \emph{involutive}\index{Yang--Baxter equation!involutive solution} 
if $r^2=\id$, and \emph{non-degenerate}\index{Yang--Baxter equation!non-degenerate solution} if the maps 
$\sigma_x\colon X\to X$ and 
$\tau_x\colon X\to X$ are bijective for all $x\in X$.

\begin{xca}
  Let $R$ be a radical ring. Prove that the map 
  \[
    r\colon R\times R\to R\times R,\quad
    r(x,y)=(-x+x\circ y,(-x+x\circ y)'\circ x\circ y),
  \]
  where $z'$ denotes the inverse of $z$ 
  with respect to the Jacobson operation, 
  is an involutive non-degenerate solution to the YBE. 
\end{xca}
\end{optional}

\section{Henriksen's theorem}

There are (non-unitary) rings with no maximal ideals. An easy example is given
in the following exercise: 

\begin{exercise}
    \leavevmode
\label{xca:Q_no_maximals}
    \begin{enumerate}
        \item Prove that the additive group of rational numbers is an abelian group
            with no maximal subgroups.
        \item Prove that the additive group $\Q$ of rationals with  
            zero multiplication (that is, $xy=0$ for all $x,y\in\Q$) is a ring with no maximal ideals. 
    \end{enumerate}
\end{exercise}

\index{Characteristic of a ring}
Recall that the \emph{characteristic of a ring} is defined as the least positive
integer $n$ such that $nx=0$ for all $x$. If no such $n$ exists, 
then we say that the ring is of \emph{characteristic zero}. 

\begin{exercise}
\label{xca:characteristic}
    % Let $R$ be a ring and $n=\lcm\{ |x|:x\in R\}$, where 
    % $|x|$ denotes the additive order of $x$. If $n<\infty$, then 
    % $R$ has characteristic $n$. If $n=\infty$, then $R$ has characteristic zero. 
    % is the least common multiple of 
    % the additive orders of the elements of $R$, then  
    Let $R$ be a non-zero ring and $p$ be a prime number. 
    If $px=0$ for all $x\in R$, then $R$ has characteristic $p$. 
\end{exercise}

We now characterize commutative rings with no maximal ideals.
We begin with some exercises. 

% Let $R$ be a commutative ring with no non-zero proper ideals. 
% Prove that if $R$ is not a field, then there exists a prime number $p$ 
% such that $R=\Z/p$ and $xy=0$ for all $x,y\in R$. 

\begin{exercise}
\label{xca:R/I_field_or_zero}
    Let $R$ be a commutative ring 
    and $I$ be an ideal of $R$. Prove that
    $I$ is maximal if and only if $R/I$ is a field
    or a ring isomorphic to 
    $\Z/p$ with zero multiplication for some prime number $p$. 
\end{exercise}

\begin{exercise}
    \label{xca:J(R)_fields}
    Let $R$ be a commutative ring. Prove that
    $J(R)$ equals the intersection of maximal ideals $M$ such that $R/M$ is a field. 
\end{exercise}

The following result appeared in~\cite{MR0424776}.

% \begin{sol}{xca:characteristic}
%     Let $n$ be the characteristic of $R$. 
%     For each $x\in R$, $x\ne 0$, $px=0$. Since the additive order $|x|$ of $x$
%     divides $p$, it follows that $|x|=p$. Since $nx=0$ for all $x$, 
%     $p$ divides $n$. The minimality of the characteristic implies that $p=n$. 
% Assume that $n<\infty$. 
% Let $X=\{n:nx=0\text{ for all $x\in R$}\}$. By definition, 
% \[ 
% m\in X
% \Longleftrightarrow |x|\text{ divides }m
% \Longleftrightarrow n\leq m
% \Longleftrightarrow n\in X.
% \]
%\end{sol}

% Let $n>0$ be the characteristic of $R$. 
% If $px=0$ for all $x$, then $n=p$. Since $px=0$ for all $x$, $p\geq n$. 
% Assume that $p>n$, then $p+t=n$ for some $n$. 
% Then $0=px=(n-t)x=nx-tx=-tx$ and hence $tx=0$ for all $x$, a contradiction. 
% If $p\nmid n$, then 
% $n=pm$. Then $px=$


\begin{theorem}[Henriksen]
\label{thm:Henriksen}
\index{Henriksen's theorem}
Let $R$ be a commutative ring. Then $R$ has
no maximal ideals if and only if 
$J(R)=R$ and $R^2+pR=R$ for every prime number $p$. 
\end{theorem}

\begin{proof}
    Assume first that $R$ has no maximal ideals. Then $J(R)=R$ by 
    Exercise~\ref{xca:J(R)_fields}. Let $p$ be a prime number
    such that $I=R^2+pR\ne R$. Then $I$ is a proper ideal of $R$. 
    Let $\pi\colon R\to R/I$ be the canonical map. Since $R^2\subseteq I$, 
    $0=\pi(xy)=\pi(x)\pi(y)$ for all $x,y\in R$. Thus $R/I$ has zero multiplication. 
    Moreover, by Exercise~\ref{xca:characteristic}, 
    $R/I$ has characteristic $p$, as $pR\subseteq I$.  
    %as $0=\pi(px)=p\pi(x)$ for all $x\in R$, 
    % since $pR\subseteq I$. 
    Thus $R/I$ is a vector space 
    over the field $\Z/p$.
    Let $\{x_\alpha:\alpha\in\Lambda\}$ be a basis
    of $R/I$. Every element $x\in R/I$ can be written uniquely
    as a finite sum of the form 
    $x=\sum \lambda_{\alpha}x_{\alpha}$ for scalars $\lambda_\alpha$. 
    Let $A$ be the ring with underlying additive group $\Z/p$ and zero multiplication. 
    For a fixed
    $\beta\in\Lambda$, the map 
    \[
    \gamma\colon R/I\to A,\quad 
    x=\sum \lambda_{\alpha}x_{\alpha}\mapsto \lambda_{\beta}
    \]
    is a surjective ring homomorphism. 
    The composition $f=\gamma\pi\colon R\to R/I\to A$ is a ring homomorphism. By Exercise~\ref{xca:R/I_field_or_zero}, $\ker f$ is a maximal ideal, a contradiction. 

    Conversely, let $M$ be a maximal ideal of $R$. If $R/M$ is a field, 
    then $J(R)\subseteq M\ne R$, a contradiction. 
    By Exercise~\ref{xca:R/I_field_or_zero},  there exists a prime
    number $p$ such that $R/M\simeq\Z/p$ as abelian groups, and the 
    multiplication in $R/M$ is zero (i.e., $xy\in M$ for all $x,y\in R$). 
    Let $\pi\colon R\to R/M$ be the canonical map. Note that 
    $R^2\subseteq M$. Moreover, $pR\subseteq M$, as 
    $\pi(px)=p\pi(x)=0$ for all $x\in R$. Thus $R^2+pR\subseteq M\ne R$, a contradiction. 
\end{proof}

We now present a non-trivial concrete example of a ring 
with no maximal ideals. For that purpose, we will use
the field of fractions $\R(X)$ of the real polynomial ring $\R[X]$.  

\begin{exercise}
    \label{xca:example}
    Let $R$ be the subring of $\R(X)$ consisting of the rational functions
    $f(X)/g(X)$, where $f(X),g(X)\in\R[X]$ and $g(0)\ne0$.
    Prove the following statements:
    \begin{enumerate}
        \item $R$ is an integral domain with a unique maximal ideal $M=XR$.
        \item $M$ has no maximal ideals. 
    \end{enumerate}  
\end{exercise}

%The construction of the previous example does not work 
%if the field of real numbers is replaced by a field of positive characteristic. 

\begin{definition}
\index{Ring!nil}
A ring $R$ is said to be \emph{nil} if for every $x\in R$ there
exists $n=n(x)\geq1$ such that $x^n=0$. 
\end{definition}

\begin{exercise}
    %Prove that a nil ring is a radical ring. 
    Prove that every nil ring is radical.
\end{exercise}

\begin{exercise}
    Let $\R[\![X]\!]$ be the ring of formal power series with real coefficients. Prove that the ideal 
    $X\R[\![X]\!]$ consisting of power series with zero constant term is a radical ring
    that is not nil. 
\end{exercise}

\begin{theorem}
	\label{thm:J(R/J)=0}
	If $R$ is a ring, then $J(R/J(R))=\{0\}$.
\end{theorem}

\begin{proof}
	If $R$ is radical, the result is trivial. Suppose then that 
	$J(R)\ne R$. Let $M$ be an arbitrary simple $R$-module. Then $M$ is 
	a simple module over $R/J(R)$ with 
	\begin{equation}
	    \label{eq:R/J(R)-module}
     	(x+J(R))\cdot m=x\cdot m,\quad
		x\in R,\,m\in M.
	\end{equation}
 
  Let us prove first that \eqref{eq:R/J(R)-module} is well-defined: If
  $r,s\in R$ are such that $r+J(R)=s+J(R)$, then $r-s\in J(R)$. Since $J(R)$
  is the intersection of the annihilators of simple $R$-modules (see
  Definition~\ref{def:J}) and $M$ is a simple $R$-module, $(r-s)\cdot
  M=\{0\}$. Thus $(r-s)\cdot m=0$ for all $m\in M$. Hence $r\cdot m=s\cdot m$
  for all $m\in M$. 

  We now prove that $M$ is a simple $R/J(R)$-module. 
  Note first that 
  \[
    (R/J(R))\cdot M=R\cdot M\ne\{0\}.
  \]
  Now if $N$ is a submodule of the $R/J(R)$-module $M$, then 
  \[
    (r+J(R))\cdot N\subseteq N
  \]
  for all $r\in R$. Thus $r\cdot N\subseteq N$ for all $r\in R$. Thus $N$ is a submodule of the $R$-module $M$. Since $M$ is simple, either $N=\{0\}$ or $N=M$. 

  Finally, if $r+J(R)\in J(R/J(R))$, then  $r\cdot M=(r+J(R))\cdot M=\{0\}$.
  Since $M$ was arbitrary, $r$ annihilates every simple $R$-module.  Thus $r\in
  J(R)$.
\end{proof}


% \begin{theorem}
% 	Let $R$ be a ring and $n\in\Z_{>0}$. Then $J(M_n(R))=M_n(J(R))$. 
% \end{theorem}

% \begin{proof}
% 	We first prove that $J(M_n(R))\subseteq M_n(J(R))$. 
% 	If $J(R)=R$, the theorem is clear. Let us assume that $J(R)\ne R$ and let  
% 	$J=J(R)$. 
% 	If $M$ is a simple $R$-module, then $M^n$ is a simple $M_n(R)$-module with the usual multiplication. 
% 	Let $x=(x_{ij})\in J(M_n(R))$ and $m_1,\dots,m_n\in M$. Then
% 	\[
% 		x\colvec{3}{m_1}{\vdots}{m_n}=0.
% 	\]
% 	In particular, $x_{ij}\in\Ann_R(M)$ for all $i,j\in\{1,\dots,n\}$. Hence 
% 	$x\in M_n(J)$. 

% 	We now prove that $M_n(J)\subseteq J(M_n(R))$. Let 
% 	\[
% 		J_1=\begin{pmatrix}
% 			J & 0 & \cdots & 0\\
% 			J & 0 & \cdots & 0\\
% 			\vdots & \vdots & \ddots & \vdots\\
% 			J & 0 & \cdots & 0
% 		\end{pmatrix}
% 		\quad\text{and}\quad
% 		x=\begin{pmatrix}
% 			x_1 & 0 & \cdots & 0\\
% 			x_2 & 0 & \cdots & 0\\
% 			\vdots & \vdots & \ddots & \vdots\\
% 			x_n & 0 & \cdots & 0
% 		\end{pmatrix}\in J_1.
% 	\]
% 	Since $x_1$ is quasi-regular, there exists $y_1\in R$ such that $x_1+y_1+x_1y_1=0$.
% 	If
% 	\[
% 		y=\begin{pmatrix}
% 			y_1 & 0 & \cdots & 0\\
% 			0 & 0 & \cdots & 0\\
% 			\vdots & \vdots & \ddots & \vdots\\
% 			0 & 0 & \cdots & 0
% 		\end{pmatrix}, 
% 	\]
% 	then $u=x+y+xy$ is lower triangular, as  
% 	\[
% 		u=\begin{pmatrix}
% 			0 & 0 & \cdots & 0\\
% 			x_2y_1 & 0 & \cdots & 0\\
% 			x_3y_1 & 0 & \cdots & 0\\
% 			\vdots & \vdots & \ddots & \vdots\\
% 			x_ny_1 & 0 & \cdots & 0
% 		\end{pmatrix}.
% 	\]
% 	Since  
% 	$u^n=0$, the element
% 	\[
% 	v=-u+u^2-u^3+\cdots+(-1)^{n-1} u^{n-1}
% 	\]
% 	is such that 
% 	$u+v+uv=0$. Thus $x$ is right quasi-regular, as  
% 	\begin{align*}
% 		x+(y+v+yv)+x(y+v+yv)&=0,
% 	\end{align*}
% 	and therefore $J_1$ is right quasi-regular. Similarly one proves that 
% 	each $J_i$ is right quasi-regular and hence $J_i\subseteq J(M_n(R))$ for all 
% 	$i\in\{1,\dots,n\}$. In conclusion, 
% 	\[
% 	J_1+\cdots+J_n\subseteq J(M_n(R))
% 	\]
% 	and therefore $M_n(J)\subseteq J(M_n(R))$.
% \end{proof}

\begin{theorem}
    Let $R$ be a ring. Then
    \[
    J(M_3(R))=M_3(J(R)).
    \]
\end{theorem}

\begin{proof}
    Let $J=J(R)$. We first prove that
    \[
    J(M_3(R))\subseteq M_3(J).
    \]
    If $J=R$, this inclusion is immediate. Otherwise, let $M$ be an
    arbitrary simple $R$-module. Then $M^3$ is a simple
    $M_3(R)$-module with the usual matrix operation. 

    Let $x=(x_{ij})\in J(M_3(R))$. Since $x$ annihilates $M^3$,
    \[
    x\colvec{3}{m_1}{m_2}{m_3}=0
    \]
    for all $m_1,m_2,m_3\in M$. By taking all coordinates except the
    $j$-th one to be zero, it follows that
    \[
    x_{ij}\cdot M=\{0\}
    \]
    for every $i,j\in\{1,2,3\}$. Since $M$ was arbitrary,
    $x_{ij}\in J(R)$ for all $i,j$. Hence
    $x\in M_3(J)$.

    We now prove that
    \[
    M_3(J)\subseteq J(M_3(R)).
    \]
    Let
    \[
    J_2=
    \begin{pmatrix}
        0 & J & 0\\
        0 & J & 0\\
        0 & J & 0
    \end{pmatrix}
    \]
    and let
    \[
    x=
    \begin{pmatrix}
        0 & x_1 & 0\\
        0 & x_2 & 0\\
        0 & x_3 & 0
    \end{pmatrix}
    \in J_2.
    \]
    Since $x_2\in J$, there exists $y_2\in R$ such that
    \[
    x_2+y_2+x_2y_2=0.
    \]
    Moreover,
    \[
    y_2=-x_2-x_2y_2\in J.
    \]
    Set
    \[
    y=
    \begin{pmatrix}
        0 & 0 & 0\\
        0 & y_2 & 0\\
        0 & 0 & 0
    \end{pmatrix}
    \in J_2
    \]
    and let
    \[
    u=x\circ y=x+y+xy.
    \]
    Then
    \[
    u=
    \begin{pmatrix}
        0 & x_1+x_1y_2 & 0\\
        0 & 0 & 0\\
        0 & x_3+x_3y_2 & 0
    \end{pmatrix},
    \]
    so $u^2=0$.

    Put
    \[
    z=y\circ(-u)=y-u-yu.
    \]
    Since $y,u\in J_2$ and $J_2$ is closed under the circle
    operation, $z\in J_2$. By associativity of the circle operation,
    \[
    x\circ z
    =x\circ\bigl(y\circ(-u)\bigr)
    =(x\circ y)\circ(-u)
    =u\circ(-u)
    =-u^2
    =0.
    \]
    Thus every element of $J_2$ has a right circle inverse in $J_2$.

    Since $(J_2,\circ)$ is a monoid in which every element has a
    right inverse, it is a group. Hence $J_2$ is quasi-regular.
    Therefore
    \[
    J_2\subseteq J(M_3(R)).
    \]

    The same argument applies to the analogous column ideals
    $J_1$ and $J_3$. Consequently,
    \[
    M_3(J)=J_1+J_2+J_3\subseteq J(M_3(R)).
    \]
\end{proof}

%\begin{theorem}
%	Let $R$ be a ring. Then $J(M_3(R))=M_3(J(R))$. 
%\end{theorem}
%
%\begin{proof}
%	Let $x=(x_{ij})\in J(M_3(R))$ and $m_1,m_2,m_3\in M$. Then
%	\[
%		x\colvec{3}{m_1}{m_2}{m_3}=0.
%	\]
%	In particular, $x_{ij}\in\Ann_R(M)$ for all $i,j\in\{1,2,3\}$. Hence 
%	$x\in M_3(J)$. 
%
%	We now prove that $M_3(J)\subseteq J(M_3(R))$. Let 
%	\[
%		J_2=\begin{pmatrix}
%			0 & J & 0\\
%			0 & J & 0\\
%			0 & J & 0
%		\end{pmatrix}
%		\quad\text{and}\quad
%		x=\begin{pmatrix}
%			0 & x_1 & 0\\
%			0 & x_2 & 0\\
%			0 & x_3 & 0
%		\end{pmatrix}\in J_2.
%	\]
%  Since $x_2$ is quasi-regular, there exists $y_2\in R$ such that
%  $x_2+y_2+x_2y_2=0$.  For 
%	\[
%		y=\begin{pmatrix}
%			0 & 0 & 0\\
%			0 & y_2 & 0\\
%			0 & 0 & 0
%		\end{pmatrix}, 
%	\]
%	let 
%    \[
%		u=x+y+xy=\begin{pmatrix}
%			0 & x_1+x_1y_2 & 0\\
%			0 & x_2+y_2+x_2y_2 & 0\\
%			0 & x_3+x_3y_2 & 0
%		\end{pmatrix}
%        =\begin{pmatrix}
%			0 & x_1+x_1y_2 & 0\\
%			0 & 0 & 0\\
%			0 & x_3+x_3y_2 & 0
%		\end{pmatrix}.
%	\]
%	Since  
%	$u^2=0$, $x$ is right quasi-regular, as  
%	\begin{align*}
%		x+(y-u-yu)+x(y-u-uy)=0,
%	\end{align*}
%  and therefore $J_2$ is right quasi-regular. Similarly one proves that both
%  $J_1$ and $J_3$ are right quasi-regular. Thus $J_i\subseteq J(M_3(R))$ for
%  all $i\in\{1,2,3\}$. In conclusion, 
%	\[
%	J_1+J_2+J_3\subseteq J(M_3(R))
%	\]
%	and therefore $M_3(J)\subseteq J(M_3(R))$.
%\end{proof}

\begin{exercise}
 	Let $R$ be a ring and $n\in\Z_{>0}$. Prove that $J(M_n(R))=M_n(J(R))$. 
\end{exercise}

\begin{exercise}
	Let $R$ be a unitary ring. Then  
	\[
	J(R)=\bigcap\{M:\text{$M$ is a maximal left ideal}\}.
	\]
\end{exercise}

\begin{exercise}
\label{xca:Jcon1}
	Let $R$ be a unitary ring and $x\in R$. The
	following statements are equivalent: 
	\begin{enumerate}
		\item $x\in J(R)$.
		\item $x\cdot M=\{0\}$ for every simple $R$-module $M$.
		\item $x\in P$ for every left primitive ideal $P$. % primitive left?
		\item $1+rx$ is invertible for all $r\in R$.
		\item $1+\sum_{i=1}^n r_ixs_i$ is invertible 
		    for all $n\geq1$ and all $r_i,s_i\in R$.
		\item $x$ belongs to every maximal left ideal. 
	\end{enumerate}
\end{exercise}

\begin{exercise}
    Let $D$ be a division ring and $R=D[X_1,\dots,X_n]$, 
    where $X_1,\ldots,X_n$ are commuting central indeterminates. 
    Prove that 
    $J(R)=\{0\}$. 
%     Como las unidades de $R$ son los elementos no nulos de $D$,
% 	$J(R)$ es un ideal de $R$. Como $D$ es simple, $J(R)\in\{0,D\}$. Si
% 	$J(R)=D$, entonces existe  $f\in R$ tal que $-1+f+(-1)f=0$ y luego $-1=0$,
% 	una contradicción. Luego $J(R)=0$.
\end{exercise}


\begin{example}
\index{Ring!local}
    A commutative and unitary ring $R$ is \emph{local} if it contains
    only one maximal ideal. 
	If $R$ is a local ring and $M$ is its maximal ideal, then $J(R)=M$. Some particular cases: 
	\begin{enumerate}
		\item If $K$ is a field and $R=K[\![X]\!]$, then $J(R)=(X)$. 
		\item If $n\geq1$, $p$ is a prime number and $R=\Z/p^n$, then $J(R)=(p)$. 
	\end{enumerate}
\end{example}

We finish the discussion on the Jacobson radical with 
some results in the case of unitary algebras. We first need an application of Zorn's lemma. 

\begin{exercise}
\label{xca:maximal_regular}
    Let $I$ be a proper left ideal that is left regular. Prove that $I$ is contained in a maximal left ideal 
    which is regular. 
\end{exercise}

% explain ideals of algebras without one. 

\begin{proposition}
	Let $A$ be a $K$-algebra and $I$ be a subset of $A$. Then $I$ is 
	a regular maximal left ideal of the algebra $A$ if and only if $I$ is 
	a regular maximal left ideal of the ring~$A$.
\end{proposition}

\begin{proof}
	Let $I$ be a regular maximal left ideal of the ring $A$. We claim that
	$\lambda I\subseteq I$ for all $\lambda\in K$. Assume that 
	$\lambda I\not\subseteq I$ for some $\lambda$. Then $I+\lambda I$
	is a left ideal of the ring $A$ that contains $I$, as 
	\[
	a(I+\lambda I)=aI+a(\lambda I)\subseteq I+\lambda (aI)\subseteq I+\lambda I.
	\]
	Since $I$ is maximal, it follows that $I+\lambda I=A$. 
	The left regularity of $I$ implies that there exists $e\in A$
	such that 
	$a-ae\in I$ for all $a\in A$. Write $e=x+\lambda y$ for $x,y\in
	I$. Then 
	\[
		e^2=e(x+\lambda y)=ex+e(\lambda y)=ex+(\lambda e)y\in I.
	\]
	Since $e-e^2\in I$ and $e^2\in I$, it follows that $e\in I$. Thus $A=I$, as
	$a-ae\in I$ for all $a\in A$, a contradiction.

	Conversely, if $I$ is a regular maximal left ideal of the algebra $A$, then 
	$I$ is a proper regular left ideal of the underlying ring $A$. We claim that $I$ is a maximal left ideal of the ring of $A$. 
    By Exercise~\ref{xca:maximal_regular}, there exists a regular maximal left ideal $M$ of the underlying ring $A$ containing $I$. By the first implication, $M$ is a regular maximal left ideal
    of the algebra $A$.
    Since $I\subseteq M\subsetneq A$ and $I$ is a maximal left ideal of the algebra $A$, it follows that $M=I$. Thus $I$ is a regular
    maximal left ideal of the underlying ring $A$.  
\end{proof}

For algebras, the Jacobson radical of an 
algebra can be defined as 
the intersection of the left ideals (of the algebra) 
that are maximal and regular. The previous 
proposition then implies that the Jacobson radical of an algebra coincides
with the Jacobson radical of the underlying ring. 

% \begin{exercise}
%     Let $A$ be an algebra. Prove that the Jacobson
%     radical of the ring $A$ coincides with the Jacobson radical of the algebra $A$. 
% \end{exercise}

% \begin{proof}
% 	Es consecuencia del teorema anterior y de que el radical de Jacobson es la
% 	intersección de los ideales a izquierda maximales y regulares.
% \end{proof}

\section{Amitsur's theorem}

We now prove an important result of Amitsur that
has several interesting applications. We first need a lemma. 

\begin{lemma}
	\label{lemma:algebraico=nil}
	Let $A$ be an algebra with one and let $x\in J(A)$. 
	Then $x$ is algebraic if and only if $x$ is nilpotent. 
\end{lemma}

\begin{proof}
    If $x$ is nilpotent, then it is algebraic. Conversely, assume that
    $x$ is algebraic.
    Then there exist $a_0,\dots,a_n\in K$ 
    not all zero such that 
    \[
		a_0+a_1x+\cdots+a_nx^n=0.
	\]
	Let $r$ be the smallest integer such that $a_r\ne 0$. Then 
	\[
		x^r(1+b_1x+\cdots+b_mx^m)=0,
	\]
	for some $b_1,\dots,b_m\in K$. 
    Since $J(A)$ is an ideal, $b_1x+\cdots+b_mx^m\in J(A)$. Hence $1+b_1x+\cdots+b_mx^m$ is a unit by
    Exercise~\ref{xca:Jcon1}, and therefore $x^r=0$.
%     Since $J(A)$ is an ideal,
%   $1+b_1x+\cdots+b_mx^m\in J(A)$ and a unit by 
% 	Exercise~\ref{xca:Jcon1}. Thus $x^r=0$.
\end{proof}

An application:

\begin{proposition}
	\label{pro:algebraica=>Jnil}
	If $A$ is an algebraic algebra with one, then $J(A)$ is the largest nil ideal of $A$.
\end{proposition}

\begin{proof}
	The previous lemma implies that $J(A)$ is a nil ideal. 
	Proposition~\ref{pro:nilJ} now implies that $J(A)$ is the largest nil ideal of $A$. 
\end{proof}

\begin{theorem}[Amitsur]
	\label{thm:Amitsur}
	\index{Amitsur's theorem}
	Let $A$ be a $K$-algebra with one such that $\dim_KA<|K|$ (as cardinals). Then 
	$J(A)$ is the largest nil ideal of $A$. 
\end{theorem}

\begin{proof}
    The result is immediate if $A=\{0\}$. Assume that $A\ne\{0\}$.

    If $K$ is finite, then $A$ is a finite-dimensional algebra. In particular, $A$ is algebraic and
	hence $J(A)$ is a nil ideal by Proposition~\ref{pro:algebraica=>Jnil}.

	Assume that $K$ is infinite and let $a\in J(A)$. Exercise~\ref{xca:Jcon1} implies that 
	every element of the form 
	$1-\lambda^{-1}a$, $\lambda\in K\setminus\{0\}$, is invertible. Thus  
	\[
		a-\lambda=-\lambda(1-\lambda^{-1}a)
	\]
  is invertible for all $\lambda\in K\setminus\{0\}$.  Let
	$S=\{(a-\lambda)^{-1}:\lambda\in K\setminus\{0\}\}$. Since 
	\[
	(a-\lambda)^{-1}=(a-\mu)^{-1}\Longleftrightarrow\lambda=\mu,
	\]
	it follows that $|S|=|K\setminus\{0\}|=|K|>\dim_KA$. Then $S$  
	is a linearly dependent set, so there are $\beta_1,\dots,\beta_n\in K$
	not all zero and distinct elements $\lambda_1,\dots,\lambda_n\in K$ such that 
	\begin{equation}
		\label{eq:Amitsur}
		\sum_{i=1}^n \beta_i(a-\lambda_i)^{-1}=0.
	\end{equation}
   Note that all the
  elements $a-\lambda_i$ and their inverses commute with one another. 
	Multiplying~\eqref{eq:Amitsur} by $\prod_{i=1}^n(a-\lambda_i)$ we get 
	\[
		\sum_{i=1}^n\beta_i\prod_{j\ne i}(a-\lambda_j)=0.
	\]
	We claim that $a$ is algebraic over $K$. Indeed,  
	\[
		f(X)=\sum_{i=1}^n\beta_i\prod_{j\ne i}(X-\lambda_j)
	\]
	is non-zero, as, for example, if $\beta_1\ne0$, then  
	$f(\lambda_1)=\beta_1(\lambda_1-\lambda_2)\cdots(\lambda_1-\lambda_n)\ne0$
	and $f(a)=0$. Since $a\in J(A)$ is algebraic, it follows that 
	$a$ is nilpotent by Lemma~\ref{lemma:algebraico=nil}. 
  Thus $J(A)$ is a nil ideal. Since every nil ideal is contained in $J(A)$ by
  Proposition~\ref{pro:nilJ}, $J(A)$ is the largest nil ideal.
\end{proof}

Amitsur's theorem implies the following result. 

\begin{corollary}
Let $K$ be an uncountable field and let $A$ be a unital
$K$-algebra of at most countable dimension. Then $J(A)$
is the largest nil ideal of $A$.
\end{corollary}

\begin{proof}
Since $\dim_K A\leq\aleph_0<|K|$, the result follows from Theorem~\ref{thm:Amitsur}.
\end{proof}

\begin{optional} 
\section{Jacobson's conjecture}

We now conclude the lecture
with two big open problems related to the Jacobson radical. The first
one is \emph{Jacobson's conjecture}. 

\begin{open}[Jacobson]
\label{prob:Jacobson}
\index{Jacobson's conjecture}
\index{Jacobson--Herstein conjecture}
Let $R$ be 
both left and right noetherian. Is then 
\[
\bigcap_{n\geq1}J(R)^n=\{0\}?
\]
\end{open}

Open problem \ref{prob:Jacobson} was originally formulated by Jacobson in 1956 \cite{MR0222106} 
for one-sided noetherian rings. In 1965 Herstein \cite{MR188253} found a counterexample
in the case of one-sided noetherian rings 
and reformulated the conjecture as it appears here. 

\begin{xca}[Herstein]
Let $D$ be the ring of rationals with odd denominators. Let
$R=\begin{pmatrix}
    D & \Q\\
    0 & \Q
\end{pmatrix}$. Prove that $R$ is right noetherian and 
$J(R)=\begin{pmatrix}
J(D) & \Q\\
0 & 0
\end{pmatrix}$. Prove that 
$J(R)^n\supseteq\begin{pmatrix}0&\Q\\0&0\end{pmatrix}$ and hence $\bigcap_nJ(R)^n$ is non-zero. 
\end{xca}
\end{optional}

\begin{optional}
\section{K\"othe's conjecture}

The following is one of the central open problems in noncommutative ring
theory.

\begin{open}[K\"othe]
\label{prob:Koethe}
\index{Kothe's conjecture@K\"othe's conjecture}
Let $R$ be a ring. Is the sum 
of two arbitrary nil left ideals of $R$ nil?
\end{open}

Open problem~\ref{prob:Koethe} known as \emph{K\"othe's conjecture}. 
The conjecture was first formulated in 1930, see \cite{MR1545158}. It is known to be true
in several cases. In full generality, the problem is still open. In~\cite{MR306251} 
Krempa proved that
the following statements are equivalent:
\begin{enumerate}
    \item K\"othe's conjecture is true.  
    \item If $R$ is a nil ring, then $R[X]$ is a radical ring. 
    \item If $R$ is a nil ring, then $M_2(R)$ is a nil ring. 
    \item Let $n\geq2$. If $R$ is a nil ring, then $M_n(R)$ is a nil ring. 
\end{enumerate}

In 1956 Amitsur formulated the following conjecture\index{Amitsur's conjecture}, see for example
\cite{MR0347873}: If $R$ is a nil ring, then $R[X]$ is a nil ring. In~\cite{MR1793911} 
Smoktunowicz found a counterexample to Amitsur's conjecture. 
This counterexample suggests that K\"othe's conjecture might be false. 
A simplification of Smoktunowicz' example
appears in~\cite{MR3169522}. See \cite{MR1879880,MR2275597} for more
information on K\"othe's conjecture and related topics. 
\end{optional}
